\documentclass[12pt]{article}
\usepackage[T1]{fontenc}
\usepackage[utf8]{inputenc}
\usepackage[brazil]{babel}
\usepackage[margin=1in]{geometry}
\usepackage{ragged2e}
\usepackage[normalem]{ulem}
\setlength{\parindent}{0pt}
\setlength{\parskip}{0.8\baselineskip}
\begin{document}
\justifying
\noindent Verifica-se que a questão objeto da controvérsia jurídica suscitada no recurso manejado pela parte recorrente esteve afetada no representativo da controvérsia - \textbf{Tema n. }\textbf{220} da TNU - (PU no processo n. PEDILEF 5004376-97.2017.4.04.7113/RS),  afetado em 22/08/2019 foi julgado (trânsito em julgado em PUIL 2266/DF (Trânsito em julgado no STJ em 25/08/2023);   e RE 1455046), por meio do qual foi elucidada a seguinte questão:

\noindent \textit{"Saber se o rol do inciso II do art. 26 c/c art. 151 da Lei nº 8.213/91 é taxativo ou se pode contemplar outras hipóteses de isenção de carência, como a gravidez de alto risco."}

\noindent Restou firmada a seguinte tese:

\noindent \textit{"1. O rol do inciso II do art. 26 da lei 8.213/91 é exaustivo. 2. A lista de doenças mencionada no inciso II, atualmente regulamentada pelo art. 151 da Lei nº 8.213/91, não é taxativa, admitindo interpretação extensiva, desde que demonstrada a especificidade e gravidade que mereçam tratamento particularizado. 3. A gravidez de alto risco, com recomendação médica de afastamento do trabalho por mais de 15 dias consecutivos, autoriza a dispensa de carência para acesso aos benefícios por incapacidade." (Tema 220 TNU).}

\noindent Com efeito, depreende-se da leitura do voto condutor do acórdão que o mesmo se encontra em conformidade com o decidido no  tema pela TNU.

\noindent A proposito, consta do acordao hostilizado: (...) Descontos Indevidos, Pedidos Genéricos Relativos aos Benefícios em Espécie, DIREITO PREVIDENCIÁRIO, Contribuição Sindical, Organização Sindical, DIREITO ADMINISTRATIVO E OUTRAS MATÉRIAS DE DIREITO PÚBLICO, Competência, Jurisdição e Competência, DIREITO PROCESSUAL CIVIL E DO TRABALHO

\noindent Nessas condições, o conteúdo do Acórdão recorrido é consentâneo com o entendimento da TNU, firmado em julgamento de recurso representativo de controvérsia, o que atrai a incidência da regra prevista pelo art. 14, III, \textit{b,} da Resolução CJF n. 586/2019:

\noindent \textit{\textbf{Art. 14.}}\textit{ Decorrido o prazo para contrarrazões, os autos serão conclusos ao magistrado responsável pelo exame preliminar de admissibilidade, que deverá, de forma sucessiva: (...) }\textit{\textbf{III }}\textit{- negar seguimento a pedido de uniformização de interpretação de lei federal interposto contra acórdão que esteja em conformidade com entendimento consolidado: }\textit{\textbf{b)}}\textit{ em }\uline{\textit{\textbf{recurso representativo de controvérsia pela Turma Nacional de Uniformização}}}\textit{ ou em pedido de uniformização de interpretação de lei dirigido ao Superior Tribunal de Justiça;(...).}

\noindent Posto isso, com arrimo no \uline{\textbf{art. 14, III, }}\uline{\textit{\textbf{b}}}\textit{,} da Resolução CJF n. 586/2019 c/c art. 1.030, I, do CPC, \textbf{NEGO SEGUIMENTO }\textbf{a}o\textbf{ PU}\textbf{.}

\noindent Intimem-se. Aguarde-se o decurso do prazo recursal. Após, certifique-se o trânsito em julgado e devolva-se ao juízo de origem.
\end{document}
